\section{Method}
\label{sec:method}

\begin{figure}[t]
\center
\includegraphics[width=\textwidth]{several_imaginations_kungfu.png}
\caption{Four imagined trajectories in \textit{KungFuMaster}. We use the same conditioning frame across the four rows, in green, and let the world model imagine the rest. As the initial frame only contains the player, there is no information about the enemies that will come next. Consequently, the world model generates different types and numbers of opponents in each simulation. It is also able to reflect an essential game mechanic, highlighted in the blue box, where the first enemy disappears after getting hit by the player.}
\vspace{0.0\linewidth}
\label{fig:kungfu}
\end{figure}

We formulate the problem as a Partially Observable Markov Decision Process (\textsc{pomdp}) with image observations $x_t \in \mathbb{R}^{h \times w \times 3}$, discrete actions $a_t \in \{1, \dots, A\}$, scalar rewards $r_t \in \mathbb{R}$, episode termination $d_t \in \{0, 1\}$, discount factor $\gamma \in (0, 1)$, initial observation distribution $\rho_0$, and environment dynamics $x_{t+1}, r_t, d_t \sim p(x_{t+1}, r_t, d_t \mid x_{\le t}, a_{\le t})$. The reinforcement learning objective is to train a policy $\pi$ that yields actions maximizing the expected sum of rewards $\mathbb{E}_\pi[\sum_{t \ge 0} \gamma^t r_t]$.

Our method relies on the three standard components to learn in imagination \citep{sutton}: experience collection, world model learning, and behavior learning. In the vein of \citet{worldmodels, simple, dreamerv1, dreamerv2}, our agent learns to act exclusively within its world model, and we only make use of real experience to learn the environment dynamics.

We repeatedly perform the three following steps: 
\begin{itemize}[left=1em]
    \setlength\itemsep{0.5em}
    \item \texttt{collect\_experience}: gather experience in the real environment with the current policy. 
    \item \texttt{update\_world\_model}: improve rewards, episode ends and next observations predictions.
    \item \texttt{update\_behavior}: in imagination, improve the policy and value functions. 
\end{itemize}

The world model is composed of a discrete autoencoder \citep{vqvae}, to convert an image to tokens and back, and a \textsc{gpt}-like autoregressive Transformer \citep{vaswani, radford_gpt2, brown_gpt3}, whose task is to capture environment dynamics. Figure \ref{fig:method} illustrates the interplay between the policy and these two components during imagination. We first describe the autoencoder and the Transformer in Sections \ref{section:im_to_tok} and \ref{section:modeling_dynamics}, respectively. Section \ref{section:learn_in_imagination} then details the procedure to learn the policy and value functions in imagination. Appendix \ref{appendix:architectures} provides a comprehensive description of model architectures and hyperparameters. Algorithm \ref{alg:iris} summarizes the training protocol.

\subsection{From image observations to tokens}
\label{section:im_to_tok}

The discrete autoencoder $(E, D)$ learns a symbolic language of its own to represent high-dimensional images as a small number of tokens. The back and forth between frames and tokens is illustrated with green arrows in Figure \ref{fig:method}.

More precisely, the encoder $E : \mathbb{R}^{h \times w \times 3} \rightarrow \{ 1, \dots, N \}^K$ converts an input image $x_t$ into $K$ tokens from a vocabulary of size $N$. Let $\mathcal{E} = \{e_i\}_{i=1}^N \in \mathbb{R}^{N \times d}$ be the corresponding embedding table of $d$-dimensional vectors. The input image $x_t$ is first passed through a Convolutional Neural Network (\textsc{cnn}) \citep{cnn_lecun} producing output $y_t \in \mathbb{R}^{K \times d}$. We then obtain the output tokens $z_t = (z_t^1, \dots, z_t^K) \in \{ 1, \dots, N \}^K$ as $z_t^k = \argmin_{i} \norm{y_t^k - e_i}_2$, the index of the closest embedding vector in $\mathcal{E}$ \citep{vqvae, vqgan}. Conversely, the \textsc{cnn} decoder $D : \{1, \dots, N\}^K \rightarrow \mathbb{R}^{h \times w \times 3}$ turns $K$ tokens back into an image. 

This discrete autoencoder is trained on previously collected frames, with an equally weighted combination of a $L_1$ reconstruction loss, a commitment loss \citep{vqvae, vqgan}, and a perceptual loss \citep{vqgan, johnson, larsen}. We use a straight-through estimator \citep{straightthrough} to enable backpropagation training. 


\subsection{Modeling dynamics}
\label{section:modeling_dynamics}

At a high level, the Transformer $G$ captures the environment dynamics by modeling the language of the discrete autoencoder over time. Its central role of unfolding imagination is highlighted with the blue arrows in Figure \ref{fig:method}. 

Specifically, $G$ operates over sequences of interleaved frame and action tokens. An input sequence $(z_0^1, \dots, z_0^K, a_0, z_1^1, \dots, z_1^K, a_1, \dots, z_t^1, \dots, z_t^K, a_t)$ is obtained from the raw sequence $(x_0, a_0, x_1, a_1, \dots, x_t, a_t)$ by encoding the frames with $E$, as described in Section \ref{section:im_to_tok}.

At each time step $t$, the Transformer models the three following distributions: 
\begin{alignat}{4}
& \text{Transition: }   &   \hat{z}_{t+1}         &   \sim p_G\big(     \hat{z}_{t+1}  &  \mid z_{\le t}, a_{\le t} \big) & \text{ with } \hat{z}_{t+1}^k \sim p_G\big(\hat{z}_{t+1}^k \mid z_{\le t}, a_{\le t}, z_{t+1}^{<k} \big)\\
& \text{Reward: }       &   \hat{r}_t \text{~~}   &   \sim p_G\big(\text{~~} \hat{r}_t &  \mid z_{\le t}, a_{\le t} \big) & \\
& \text{Termination: }  &   \hat{d}_t \text{~~}   &   \sim p_G\big(\text{~~} \hat{d}_t &  \mid z_{\le t}, a_{\le t} \big) & 
\end{alignat}

Note that the conditioning for the $k$-th token also includes $z_{t+1}^{<k} \coloneqq (z_{t+1}^1, \dots, z_{t+1}^{k-1})$, the tokens that were already predicted, i.e. the autoregressive process happens at the token level.

We train $G$ in a self-supervised manner on segments of $L$ time steps, sampled from past experience. We use a cross-entropy loss for the transition and termination predictors, and a mean-squared error loss or a cross-entropy loss for the reward predictor, depending on the reward function.

\subsection{Learning in imagination}
\label{section:learn_in_imagination}

\begin{figure}
\center
\includegraphics[width=\textwidth]{perfect_wm_pong.png}
\caption{Pixel perfect predictions in \textit{Pong}. The top row displays a test trajectory collected in the real environment. The bottom row depicts the reenactment of that trajectory inside the world model. More precisely, we condition the world model with the first two frames of the true sequence, in green. We then sequentially feed it the true actions and let it imagine the subsequent frames. After only 120 games of training, the world model perfectly simulates the ball's trajectory and players' movements. Notably, it also captures the game mechanic of updating the scoreboard after winning an exchange, as shown in the blue box.}
\label{fig:pong}
\end{figure}


Together, the discrete autoencoder $(E,D)$ and the Transformer $G$ form a world model, capable of imagination. The policy $\pi$, depicted with purple arrows in Figure \ref{fig:method}, exclusively learns in this imagination \textsc{mdp}. 

At time step $t$, the policy observes a reconstructed image observation $\hat{x}_t$ and samples action $a_t \sim \pi(a_t | \hat{x}_{\le t})$. The world model then predicts the reward $\hat{r}_t$, the episode end $\hat{d}_t$, and the next observation $\hat{x}_{t+1} = D(\hat{z}_{t+1})$, with $\hat{z}_{t+1} \sim p_G(\hat{z}_{t+1} \mid z_0, a_0, \hat{z}_1, a_1, \dots, \hat{z}_t, a_t)$. This imagination procedure is initialized with a real observation $x_0$ sampled from past experience, and is rolled out for $H$ steps, the imagination horizon hyperparameter. We stop if an episode end is predicted before reaching the horizon. Figure \ref{fig:method} illustrates the imagination procedure.

As we roll out imagination for a fixed number of steps, we cannot simply use a Monte Carlo estimate for the expected return. Hence, to bootstrap the rewards that the agent would get beyond a given time step, we have a value network $V$ that estimates $V(\hat{x}_t) \simeq \mathbb{E}_\pi\big[ \sum_{\tau \ge t} \gamma^{\tau - t} \hat{r}_\tau \big]$.

Many actor-critic methods could be employed to train $\pi$ and $V$ in imagination \citep{sutton, simple, dreamerv1}. For the sake of simplicity, we opt for the learning objectives and hyperparameters of DreamerV2 \citep{dreamerv2}, that delivered strong performance in Atari games. Appendix \ref{appendix:actor_critic} gives a detailed breakdown of the reinforcement learning objectives.